% LTeX: language=cs-CZ
\documentclass[12pt]{article}
\usepackage[czech]{babel}
\usepackage{hyphsubst}
\usepackage{graphicx}
\usepackage{indentfirst}
\usepackage{microtype}
\usepackage{float}
\usepackage{amsmath}
\usepackage{esvect}

\selectlanguage{czech}
\PassOptionsToPackage{hyphens}{url}

\usepackage{hyperref}
\hypersetup{
    colorlinks=true,
    linkcolor=black,
    filecolor=magenta,      
    urlcolor=blue,
    pdftitle={Elektroviz},
}
\urlstyle{same}

% define helper command for typesetting code
\newcommand{\code}[1]{\texttt{#1}}
\def\CC{{C\nolinebreak[4]\hspace{-.05em}\raisebox{.4ex}{\tiny\bf ++}}}

\chardef\_=`_

% increase spacing between table rows
\renewcommand{\arraystretch}{1.25}

\begin{document}

\begin{figure}[H]
	\centering
	\includegraphics[width=0.8\textwidth]{pic/LOGO_KIV_CMYK.png}
\end{figure}

\begin{center}
	\vspace{.5cm}
	\LARGE{Elektroviz - Vizualizace elektrického pole}\\
	\large{Semestrální práce - KIV/UPG 2024/25}
\end{center}

\vfill

\noindent
Západočeská univerzita v Plzni \hfill Pavel Altmann - A23B0264P\\
Katedra informatiky a výpočetní techniky \hfill 1. odevzdání\\
Stráveno hodin: 80 \hfill \today
\thispagestyle{empty}

\newpage
\setcounter{page}{1}

\tableofcontents

\newpage

\section{Zadání}

Úkolem semestrální práce je vytvoření interaktivního programu pro vizualizaci
elektrostatického pole vytvořeného bodovými elektrickými náboji rozmístěných v
prostoru.

Program bude možné spustit z příkazové řádky skriptem \code{./run.sh} s
volitelným celočíselným parametrem určujícím scénář (viz níže) k vizualizaci.
Nebude-li parametr zadán, bude zvolen scénář 0. Po spuštění programu se zobrazí
okno o minimální počáteční velikosti 800x600px, ve kterém se zobrazí přehledně
všechny bodové náboje zvoleného scénáře. Vizuální reprezentace náboje (v
nejjednodušší podobě kružnice s popiskem) musí být volena tak, aby bylo možné
jednoznačně určit velikost náboje a jeho znaménko. V místě sondy bude zobrazen
šipkou vektor intenzity vektorového pole 𝐸⃗ a opatřen popiskem udávající jeho
velikost $\|E\|$. Všechny elementy vizualizace (náboje, popisky, šipky) musí být
plně viditelné.

Intenzita vektorového pole v daném místě $\vec{x}$ se spočte
z Coulombova zákona jako:
\[ \vec{E}(\vec{x}) = \frac{1}{4 \pi \epsilon_0}
\sum_i q_i \frac{\vv{x_i - x}}{\|\vv{x_i - x}\|^3} \]

Vizualizace se bude pravidelně (s časem) aktualizovat. Jednotlivé elementy
vizualizace budou automaticky svou velikost a rozmístění přizpůsobovat velikosti
(a tvaru) okna tak, aby se prostor okna „maximálně“ využil, ale zároveň aby
nedošlo ke zkreslení souřadnic.

Scénáře:

\begin{itemize}
  \item Scénář 0 = pouze jeden bodový náboj +1C v počátku souřadnic (0, 0);
  \item Scénář 1 = jeden náboj +1C v (-1, 0), druhý +1C v (1, 0);
  \item Scénář 2 = jeden náboj -1C v (-1, 0), druhý +2C v (1, 0);
  \item Scénář 3 = čtyři náboje: +1C v (-1, -1), +2C v (1, -1), -3C v (1, 1)
    a -4C v (-1, 1);
\end{itemize}

Ve všech těchto scénářích se sonda pohybuje po kružnici se středem v počátku
souřadnic (0, 0), poloměrem 1 a úhlovou rychlostí \( \frac{\pi}{6}
\frac{rad}{s} \).

\section{Implementace řešení}

\subsection{Použité technologie}

Vizualizace je naprogrogramována v jazyce \CC{} 23 s využitím 2D grafické
knihovny \href{https://www.raylib.com/}{Raylib} verze 5.5.


\subsection{Architektura}

Program využívá OOP architekturu, kde každý grafický prvek je uzavřený do své
třídy a zodpovědný za veškeré výpočty a kreslení, které se ho týkají. Funkce
main má na starost pouze počáteční nastavení a následné předávání událostí
(nový snímek, změna velikosti okna, apod.) objektům. Některé objekty jako
součást obsahují další podobjekty a potom mají na starost předávání událostí.

\subsection{Objekty}

Každý objekt má především 2 metody \code{update} a \code{draw}. \code{update} je
zodpovědná za přepočítání vlastností objektu v závislosti na čase, popř. změně
vlastností jiných objektů (např. sonda mění polohu podle času a naměřenou
hodnotu podle síly a polohy nábojů).

Každý zde popsaný objekt se nachází ve svém vlastním souboru

\subsubsection{\code{Charge}}

Reprezentuje v simulaci jeden náboj. 

Velikost náboje může být buď konstantní nebo proměnlivá v čase - reprezentováno
třídami \code{StaticStrength} a \code{VariableStrength}. 

Vykresluje se jako kruh, jehož plocha je přímo úměrná velikosti náboje. Dále má
bílé ohraničení a barvu výplně reprezentující velikost náboje - červená 5 C a více,
modrou 5 C a méně a lineární interpolaci přes tmavě šedou pro hodnoty mezi tím.

\subsubsection{\code{Probe}}

Reprezentuje v simulaci elektrickou sondu.

Může být buď otáčivě pohyblivá nebo statická (využití v třídě \code{Grid}) -
reprezentováno třídami \code{Static}, \code{Rotating} ze jmenného prostředí
\code{position}.

Vykresluje se jako šipka reprezentující směr a sílu el. pole v daném bodě.
Délka šipky odpovídá binárnímu logaritmu normy vektoru el. pole namísto přímé
úměry, neboť ta se pohybuje v rozmezí několika řádů a tedy je buď velmi malá
nebo větší než celá obrazovka. Zárověň je volitelně může zobrazovat text s
numerickou velikostí intenzity ve vědecké notaci.

\subsubsection{\code{Grid}}

Představuje mřížku na pozadí, včetně sond v jejích průsečících.

Má nezávisle nastavitelné rozestupy mezi čárami v obou směrech v pixelech - ve
výchozím nastavení to je 50x50. Je dobré nenastavovat příliš malé rozestupy,
protože pak je jednak obraz nepřehledný a zároveň se velmi zvyšují nároky na
výpočetní výkon.

Pro zrychlení kreslení se čáry i sondy generují pouze při inicializaci a při
změně rozměrů a jinak jsou uložené předpočítané v soukromých atributech třídy.

\subsection{Ostatní soubory}

\subsubsection{\code{field}}

Jmenné prostředí obsahuje funkce \code{E} a \code{potential}, počítající vektor,
resp. potenciál elektrického pole v zadaném bodě.

\subsubsection{\code{utils}}

Obsahuje pomocné funkce pro míchání barev \code{lerpColor} a \code{lerpColor3},
funkci \code{sigmoid} pro počítání stejnjmenné matematické funkce a funkci
\code{trim} pro odstranění bílých znaků z řetězce.

\subsubsection{\code{MathEval}}

Sada funkcí pro vyhodnocovaní libovolných matematických výrazů, využívaná pro
dynamicky nastavitelné vlastnosti. Další info patří do pozdějšího odevzdaní.

\section{Ovládání programu}

\subsection{Sestavení}

Program je kompatibilní s GNU/Linux systémy s podporou grafického protokolu X11.

Sestavení se provádí pomocí přiloženého \code{Makefile} příkazem \code{make
setup build}.

Protože ale mnoho systémů nemá nainstalované hlavičkové soubory pro knihovnu X11
(potřebné pro vývoj většiny grafických aplikací), má zastaralé verze překladačů
(obzvlástě ty založené na Debianu), apod. zvolil jsem použití systému Docker.
Sestavovací skript \code{./build.sh} vytvoří kontejner z přiloženého souboru
\code{Dockerfile} a spustí ho. Ten spustí přiložený \code{Makefile}, který do
složky \code{vendor} stáhne knihovny a vše sestaví do složky \code{build}.
Skript pak ještě přesune výsledný spustitelný soubor do souboru
\code{bin/elektroviz.new}.

\subsection{Spuštění}

Program přebírá z příkazové řádky parametr \code{scenario}. a volitelně
\code{spacing}. \code{scenario} Určuje který scénář má být spuštěn po zapnutí
aplikace, při absenci parametru se použije scénář 0. a \code{spacing} určuje
rozestupy mřízky na pozadí.

\textbf{Důležité}: Program je nutné spustit ve stejné složce, v jaké se nechází
složka \code{scenarios}, jinak nebude schopný načíst žadný z výchozích scénářů.
Originálně je v \code{src} ale \code{./build.sh} jí pro snadnější manipulaci 
nakopíruje i do kořenové složky.

\section{Limitace}

Program pouze minimálně ošetřuje nesprávné uživatelské vstupy a reaguje na ně
vyjímkami, které většinou způsobují pád aplikace. 

\end{document}
